\par\noindent\textbf{(1)} 两个法向Fresnel振幅r01=(n0−n1)/(n0+n1)、r12=(n1−n2)/(n1+n2)。膜内往返相位2δ=4πn1d/λ，多次反射几何级数给r\_tot=(r01+r12e\textasciicircum{}\{2iδ\})/(1+r01r12e\textasciicircum{}\{2iδ\})，对无吸收膜成立。
\par\noindent\textbf{(2)} 若光学厚度n1d=λ/4，则δ=π/2、e\textasciicircum{}\{2iδ\}=−1，总反射简化为r\_tot=(r01−r12)/(1−r01r12)。两界面反射相位差π，若两反射幅大小匹配便可完全相消。
\par\noindent\textbf{(3)} 零反射要求r01=r12，即(n0−n1)/(n0+n1)=(n1−n2)/(n1+n2)，交叉相乘化简得n1²=n0n2，因此最优无损膜折射率n1=√(n0n2)，厚度d=λ/(4n1)；这不是无条件任意膜都可达到的消反射。
\par\noindent\textbf{(4)} 偏离设计波长会改变δ，反射率回升；吸收介质n1复数时反射零点需同时匹配实虚阻抗，简单几何平均不够。若n0=n2，最优n1=n0意味着无界面反射，任意d均无反射，四分之一波厚已不再必要。
